\section{Cel i założenia eksperymentu}
Głównym celem eksperymentu jest ocena skuteczności wybranych metod obrony dużych modeli językowych przed bezpośrednimi ataki promptowymi. Skupiono się na metodach stosowanych na etapie wdrożenia modelu.
Wybrano trzy podejścia do mitygacji tej podatności, które różnią się od siebie miejscem integracji z modelem.
Jedną z nich jest metoda obrony przy użyciu inżynierii aktywacji, inspirowana badaniami Panickssery~i~in.~\cite{panickssery2024steeringllama2contrastive} oraz Wang~i~in.~\cite{wang2024inferalignerinferencetimealignmentharmlessness}.
Następną sprawdzoną techniką jest obrona na poziomie promptu przy użyciu metody Self-Reminder~\cite{wu2023selfreminder}.
Ostatnim wybranym mechanizmem jest filtrowanie zewnętrzne przy użyciu modelu Llama-Guard~\cite{inan2023llamaguardllmbasedinputoutput} do analizy zapytań wejściowych oraz odpowiedzi wyjściowych.
Skuteczność metod obrony poddano ocenie przy użyciu trzech odmiennych strategii ataku.
Pierwszą z nich jest białoskrzynkowy atak GCG~\cite{zou2023universaltransferableadversarialattacks}, który oparty jest na gradientowej optymalizacji sufiksu tokenów dodawanych do szkodliwego promptu.
Drugą jest czarnoskrzynkowy, automatyczny atak PAIR~\cite{chao2024jailbreakingblackboxlarge}, którego celem jest zautomatyzowane generowanie niebezpiecznych zapytań, wykorzystując do tego inne LLM.
Ostatnim podejściem do przełamania zabezpieczeń modelu jest manualna metoda ataku DAN~\cite{choi2025dan} opartej o odgrywanie ról, w której narzuca się modelowi osobowość cechującą się brakiem ograniczeń.
Dodatkowo, przeanalizowano kompromis między bezpieczeństwem a użytecznością modelu oraz sprawdzono wpływ wprowadzonych metod obrony na opóźnienia odpowiedzi modelu.